
Machine learning's reproducibility crisis persists despite reproducibility badges. Major conferences including NeurIPS, ICML, and ICLR now incentivize artifact submission through badge programs, requiring authors to deposit code, data cards, and documentation. These programs have proliferated since 2018, creating a substantial corpus of artifacts across thousands of published papers. Yet their effectiveness remains empirically unverified---do these artifacts actually improve reproducibility?

The crisis is well-documented. Kapoor and Narayanan (2023) identified data leakage affecting 294 papers across 17 fields, establishing an 8-type taxonomy of methodological errors. Semmelrock et al. (2024) developed a comprehensive framework cataloging reproducibility barriers, mapping technical and procedural obstacles that prevent independent verification of published results. These studies establish that reproducibility problems exist at scale, but focus on identifying barriers rather than measuring the effectiveness of interventions designed to address them.

Reproducibility badges represent a policy-level intervention: by requiring artifact deposition, venues aim to increase the availability of implementation details that enable reproduction. However, this intervention operates on an untested assumption---that artifact presence improves reproducibility outcomes. No empirical study has quantified whether badges succeed at this goal. The gap is critical: if badges create checkbox compliance without substantive quality, resources are wasted and the reproducibility crisis persists despite apparent progress.

The deeper challenge is distinguishing artifact presence from artifact quality. Documentation frameworks exist---FAIR principles for scientific data \citep{Wilkinson2016}, Croissant-RAI metadata standards \citep{Jain2024}, and venue-specific badge requirements---but they focus on what artifacts should contain, not whether deposited artifacts actually meet these standards. Recent evidence suggests compliance rates are low: Gim et al. (2025) found only 5\% of medical imaging datasets meet FAIR ''Findable'' criteria and 0\% meet ''Reusable'' criteria. The machine learning community lacks quantitative measurement linking artifact quality to reproducibility outcomes.

Our key insight enables large-scale measurement: performance variance across independent reproduction attempts provides a scalable proxy for reproducibility. Traditional reproducibility studies require manual replication by independent teams (e.g., NeurIPS Reproducibility Challenge), limiting sample sizes to tens of papers. In contrast, Papers with Code aggregates performance results from thousands of independent reproduction attempts across 4,000+ benchmarks. When multiple independent research groups report similar performance on a benchmark, the evaluation protocol exhibits high procedural consistency---a necessary condition for reproducibility. Conversely, high variance suggests ambiguous specifications that different groups interpret inconsistently.

We operationalize this insight by measuring the coefficient of variation (CV = $\sigma /\mu )$ of reported performance across independent reproduction attempts, treating lower CV as an indicator of higher reproducibility. This enables quantitative analysis at scale of the relationship between documentation artifact quality and reproducibility outcomes---the first continuous quality-outcome measurement in ML benchmarks.

We measure artifact quality across 108 classification benchmarks published 2019-2024 with $\geq 5$ independent reproduction attempts each. Using a validated rubric covering preprocessing specifications, data split protocols, evaluation procedures, and hyperparameters, we find mean artifact quality of 2.43/10 (threshold: 7.0), with critical gaps in evaluation protocols (1.19/10) and hyperparameters (1.16/10). Automated scoring consistency $\kappa$=1.0 confirms measurement reliability. Despite artifact deposition at scale, we observe no statistically significant reduction in performance variance: benchmarks with $\geq 2$ artifacts (GitHub repo, dataset card, badge) show mean CV=0.035 compared to CV=0.069 for benchmarks with fewer artifacts (Mann-Whitney p=0.418, Cohen's d=0.464). While the directional trend is positive, the effect size falls below the medium threshold (d=0.5) and statistical power is limited by sample size (n=22 benchmarks in final analysis, vs. target n=100).

This work makes three contributions. First, we provide the first continuous quality-outcome measurement linking artifact quality to reproducibility in ML benchmarks, revealing that artifact presence does not guarantee artifact quality. Where prior work used binary FAIR compliance \citep{Gim2025}, we introduce continuous quality scoring (0-10) linked to variance outcomes. Reproducibility badges succeed at increasing artifact deposition but fail to enforce quality standards. Second, we operationalize performance variance (CV) as a scalable reproducibility proxy, demonstrating its viability for meta-analyses that would be infeasible with manual replication studies. Third, we report a carefully characterized null result with positive directional trend, providing effect size estimates (d=0.464) and power analysis for future work. Our findings suggest that badge programs require quality enforcement mechanisms, not merely presence incentives, to meaningfully impact reproducibility outcomes.
